\textbf{Keywords:} Checklist, Specification, Design, Tapeout, Schematic
Rules, Layout Rules

\section{Checklist}\label{checklist}

There are roughly 3 phases of analog design.

\begin{itemize}
\tightlist
\item
  Specification
\item
  Design
\item
  Tapeout
\end{itemize}

The specification phase is where you think deeply through the design.

Can the design meet the key parameters you need? How will I verify the
circuit? Do I know how to make the circuit?

These questions and more, are so common that most companies will have
checklists that we use when we review the specification, design, and
tapeout.

These checklists are closely guarded secrets, as the content contain
significant amount of knowledge accumulated over numerous blunders,
mistakes, failure to imagine, and physics teaching us a lesson.

The design phase is where we make the schematic, and simulate the
schematic. We explore circuit architectures, fix problem corners, check
our design over temperature, voltage, process corners (slow transistors,
fast transistors, mismatch)

The tapeout phase is where we translate the schematic into layout, check
the design rules (DRC), do layout versus schematic (LVS) and extract
circuit parasitics to check layout parasitic effects (LPE). And, of
course, simulate most things again.

I've made a checklist below for the most common questions that you need
to ask yourself.

\subsection{Specification}\label{specification}\index{specification@Specification}

\begin{longtable}[]{@{}
  >{\raggedright\arraybackslash}p{(\linewidth - 6\tabcolsep) * \real{0.1364}}
  >{\raggedright\arraybackslash}p{(\linewidth - 6\tabcolsep) * \real{0.7898}}
  >{\raggedright\arraybackslash}p{(\linewidth - 6\tabcolsep) * \real{0.0284}}
  >{\raggedright\arraybackslash}p{(\linewidth - 6\tabcolsep) * \real{0.0455}}@{}}
\toprule\noalign{}
\begin{minipage}[b]{\linewidth}\raggedright
Item
\end{minipage} & \begin{minipage}[b]{\linewidth}\raggedright
Description
\end{minipage} & \begin{minipage}[b]{\linewidth}\raggedright
Yes
\end{minipage} & \begin{minipage}[b]{\linewidth}\raggedright
Action
\end{minipage} \\
\midrule\noalign{}
\endhead
\bottomrule\noalign{}
\endlastfoot
Functional description & Have you described what the IP shall do? & & \\
Key parameters & Have you updated your key parameters in the README &
& \\
Architecture & Have you described the circuit architecture? How should
it work? & & \\
Realism & Do you know how to do what you plan to do? & & \\
Verification plan & Have you described exactly what you need to check?
For example, stability of OTAs, current consumption, key parameters &
& \\
Specification & Have you added a specification for all parameters you
intend to check. For example, phase margin should always be larger than
45 degrees. & & \\
\end{longtable}

\subsection{Design}\label{design}\index{design@Design}

\begin{longtable}[]{@{}
  >{\raggedright\arraybackslash}p{(\linewidth - 6\tabcolsep) * \real{0.1608}}
  >{\raggedright\arraybackslash}p{(\linewidth - 6\tabcolsep) * \real{0.7483}}
  >{\raggedright\arraybackslash}p{(\linewidth - 6\tabcolsep) * \real{0.0350}}
  >{\raggedright\arraybackslash}p{(\linewidth - 6\tabcolsep) * \real{0.0559}}@{}}
\toprule\noalign{}
\begin{minipage}[b]{\linewidth}\raggedright
Item
\end{minipage} & \begin{minipage}[b]{\linewidth}\raggedright
Description
\end{minipage} & \begin{minipage}[b]{\linewidth}\raggedright
Yes
\end{minipage} & \begin{minipage}[b]{\linewidth}\raggedright
Action
\end{minipage} \\
\midrule\noalign{}
\endhead
\bottomrule\noalign{}
\endlastfoot
git & Have you committed the schematics to the repository? & & \\
git push & Are the schematics pushed to github? Are you sure? & & \\
git tag & Is the current version tagged & & \\
Implementation & Are all schematics described with their own markdown
file? & & \\
Verification plan & Are all items on the verification plan completed? If
not, have you described why it's no longer relevant? & & \\
Electrical parameters & Are the electrical parameters updated with
simulated results? & & \\
Spec violations & Have you explained why the specification violations
are not an issue? & & \\
Simulation & Are the required corners run (typical, slow, fast, mc) &
& \\
\end{longtable}

\subsection{Tapeout}\label{tapeout}\index{tapeout@Tapeout}

\begin{longtable}[]{@{}
  >{\raggedright\arraybackslash}p{(\linewidth - 6\tabcolsep) * \real{0.1121}}
  >{\raggedright\arraybackslash}p{(\linewidth - 6\tabcolsep) * \real{0.7664}}
  >{\raggedright\arraybackslash}p{(\linewidth - 6\tabcolsep) * \real{0.0467}}
  >{\raggedright\arraybackslash}p{(\linewidth - 6\tabcolsep) * \real{0.0748}}@{}}
\toprule\noalign{}
\begin{minipage}[b]{\linewidth}\raggedright
Item
\end{minipage} & \begin{minipage}[b]{\linewidth}\raggedright
Description
\end{minipage} & \begin{minipage}[b]{\linewidth}\raggedright
Yes
\end{minipage} & \begin{minipage}[b]{\linewidth}\raggedright
Action
\end{minipage} \\
\midrule\noalign{}
\endhead
\bottomrule\noalign{}
\endlastfoot
git & Are all schematics and layout committed? Are you sure? & & \\
git push & Have you pushed to github? & & \\
git tag & Is the latest version tagged? & & \\
LVS & Is the LVS on github passing (green)? & & \\
DRC & Is the DRC on github passing (green)? & & \\
LPE & Is the LPE on github passing (green)? & & \\
Simulation & Are the required corners re-run with layout parasitics
(typical, slow, fast, mc) & & \\
\end{longtable}

\section{Schematic rules}\label{schematic-rules}\index{schematic rules@Schematic rules}

Rules are nice. They reduce the cognitive load since a decision already
has been made. You may disagree with the rule, but in analog design it's
not that important what the ``rule'' is, but sometimes more important
that the rule is followed.

Naming rules are one example. We can have a philosophical discussion
till the end of time of whether it's best with uppercase or lower case.
CamelCase, however, is wrong in schematics and SPICE, since SPICE is
case-insensitive, so ``vref'' is the same as ``Vref''.

Below are the rules I like. You may disagree, but if you're my student,
then you don't have a choice. Follow them, or the grade may suffer.

\subsection{Only uppercase names
allowed}\label{only-uppercase-names-allowed}

\textbf{Do:} AVDD, \textbf{Don't:} aVdd

Although editors can handle mix of upper case and lower case, SPICE,
cannot. SPICE is case insensitive. That means AVDD == aVdd in SPICE, but
AVDD != aVdd in editors. A such mixing case is a bad idea, so one must
be picked, and uppercase the chosen one. Why? Why not?

\subsubsection{Use same net throughout
hierarchy}\label{use-same-net-throughout-hierarchy}

\textbf{Do:} VDD -\textgreater{} VDD -\textgreater{} VDD,
\textbf{Don't:} VDD -\textgreater{} LOCALVDD -\textgreater{} CELLVDD

Debugging becomes a lot simpler if a net keeps it's name throughout the
schematic hierarchy. Especially bad are cases where a net name is reused
on multiple levels of hierarchy. Imagine the following scenario

\begin{verbatim}
     | (sub block)          |
     |                      |
VDD -|-LVDD---/ ---- VDD    |
     |                      |
\end{verbatim}

Here the net name VDD is used on the top level, while LVDD is used in
the sub-block. In the sub-block there is a power switch between LVDD and
VDD. In this case VDD != VDD on top level, which can lead to long
debugging times. There are, however, a few exceptions. For example,
using VDD on standard cells (inverters, ANDs etc) is ok, even though the
power supply is not called VDD

\subsubsection{Spend time on making schematics
pretty}\label{spend-time-on-making-schematics-pretty}

It matters how schematics look. Think of it like this. In 10 years, you
will be asked to port, re-simulate, fix a bug, on your design. If you
have spent some time adding comments, making things look pretty, etc,
then the job will be much, much easier. ``A pretty schematic is a love
letter to your future self''.

\begin{quote}
PS: This applies to documentation, code, and everything you make.
\end{quote}

\subsubsection{All digital nets must be ``active high''
naming}\label{all-digital-nets-must-be-active-high-naming}

\textbf{Example:} PWRUP\_3V3, PWRUP\_N\_3V3

The PWRUP\_3V3 should be understood as ``When PWRUP\_3V3 is high, then
the block is powered up''

The PWRUP\_N\_3V3 should be understood as ``When PWRUP\_N\_3V3 is high
then the block is powered down''

\subsubsection{Bias currents shall be named IB\textless device sending
current
(P\textbar N)\textgreater{}}\label{bias-currents-shall-be-named-ibdevice-sending-current-pn}

\textbf{Do:} IBP\_1U, \textbf{Don't:} IBIAS\_1U

The ``P'' post-fix tells us the current comes from a PMOS, so we can put
it into a NMOS diode connected transistor. On the IBIAS we have no idea
which way the current flows.

\section{Layout rules}\label{layout-rules}\index{layout rules@Layout rules}

\subsection{Limit the amount of transistor Width's and Length's that you
use}\label{limit-the-amount-of-transistor-widths-and-lengths-that-you-use}

\textbf{Do:} W = 1.0 um, L = 180 nm, Use multiplier for other sizes

\textbf{Don't:} W1 = 1 um, W2 = 1.1 um, W3 = 1.2 um, W4 = 2 um, W5 = 2.1
um

You want the layout to be relatively regular. A bunch of different Ls
and Ws is a pain when you do layout

\subsubsection{Use pre-defined transistors for regular
layout}\label{use-pre-defined-transistors-for-regular-layout}

\textbf{Do:} Use JNW\_ATR\_SKY130A and JNW\_TR\_SKY130A

\subsubsection{Always use two fingers for analog
transistors}\label{always-use-two-fingers-for-analog-transistors}

That way, you don't have to worry about current direction in the layout.

\subsubsection{Always run gates in the same
direction}\label{always-run-gates-in-the-same-direction}

Mobility of transistors (especially PMOS) is affected by strain, so if
you rotate a transistor it will not have the same current, and change in
current as a function of stress.

I've seen ICs have to be taped out again due to rotated transistors.

\subsubsection{Always have dummy poly
gates}\label{always-have-dummy-poly-gates}

For large lengths (\textgreater{} 500 nm) the lithography effects are
not that severe, but the etching of the gate material will be asymmetric
if there are no poly dummies. Make sure the poly dummy is exactly the
same spacing on both sides.

For small lengths (\textless{} 200 nm) the lithography effects start to
matter. The light used in most litho is 193 nm. 193 nm is used all the
way down to about 7 nm.

Due to diffraction effects, it's common to have extremely regular poly
spacing, an exact distance such that the interference from neighboring
poly's align perfectly to the next poly.

\subsubsection{Always place transistors away from well
edge}\label{always-place-transistors-away-from-well-edge}

Close to the N-well edge the donor consentration will be higher. Ion
implantation is used for the wells, and the ions will scatter of the
oxide wall, and increase the doping concentration close to the edge. As
such, the transistor threshold voltage will increase close to a well
edge.

Keep transistors about 3 um away from the N-well edge if threshold
voltage is important.

\section{Would you like to know
more?}\label{would-you-like-to-know-more}

Johns and Martin, the textbook path through analog design {[}1{]}

Razavi, a second voice on the same material {[}2{]}

The JSSC and ISSCC archives, for how it is actually being done now

\phantomsection\label{refs}
\begin{CSLReferences}{0}{0}
\bibitem[\citeproctext]{ref-johns}
\CSLLeftMargin{{[}1{]} }%
\CSLRightInline{D. Johns and K. Martin, \emph{Analog integrated circuit
design}. John Wiley \& Sons, Inc., 1997. }

\bibitem[\citeproctext]{ref-razavi}
\CSLLeftMargin{{[}2{]} }%
\CSLRightInline{B. Razavi, \emph{Design of analog {CMOS} integrated
circuits}. McGraw-Hill, 2001. }

\end{CSLReferences}
